\subsection{What produces the Gram estimates}
\label{sec:actual-gram}

The following reconstruction goes below the abstract frame argument.
Its conclusions remain conditional on the explicitly identified
spectral and arithmetic inputs below.  In particular, the estimates
are not merely bounds for functions at infinity: finite-place
density and mass identities are essential to their normalization.
Source identifiers AG1--AG10 refer to the accompanying provenance ledger.

\paragraph{The kernel and its four parts (AG1--AG3).}
Write \(P_u:\mathcal V_u\to\mathcal H\) for the full period map,
\(n_u=\dim\mathcal V_u\), and
\[
 c_u=\frac{\operatorname{vol}N_u\operatorname{vol}H_u}
                 {\dim\pi_u}>0,\qquad
 A_u=c_u^{-1/2}P_u,\qquad K_{uv}=A_u^*A_v .
\]
These are the full source spaces, before the column restriction.
The absolutely convergent geometric expansion separates as
\[
 K_{uv}=\mathbf1_{u=v}I+Z_{uv}+N_{uv}+G_{uv}.                 \tag{AG.1}
\]
Here \(Z\) consists of the nonidentity terms with zero Gram
determinant; \(N\) and \(G\) are respectively the nongeneric and
generic terms with nonzero determinant.  The identity-cell
normalization gives the displayed \(I\).  For \(u\ne v\), the
zero-determinant amplitude intertwines distinct compact types
(after an inner conjugation), so Schur's lemma gives \(Z_{uv}=0\).
The geometry establishing this intertwining, and the absolute
convergence needed for (AG.1), are retained inputs.

The generic squared norm is controlled by a scalar pair sum and a
small remainder.  Set
\[
 X_{uv}=\frac{n_un_v}{D_*},\quad
 F=\prod_{\mathfrak p\ {\rm typed}}(1+3/N\mathfrak p),\quad
 E=e^{-c_Lkf\log k}.
\]
The source's pair-trace identity gives
\[
 \|G_{uv}\|_{\rm HS}^2\le FE X_{uv}+S_{uv},                 \tag{AG.2}
\]
where \(S_{uv}=F\gamma_G\operatorname{Re}\sum_{\rm scalar}p\).
The omitted pairs comprise defect \(s_0\ge2\), and the nonscalar
generic pairs of defect \(s_0=1\).  Here \(s_0=k-\dim_E W_p\)
for the matching subspace attached to the pair.  Their arithmetic counting bound is the
input \(E X_{uv}\); it is not a consequence of Schur orthogonality.

\paragraph{The actual archimedean functions (AG3--AG5).}
At a complex-place representative \(\nu\), the Bessel amplitude is
the double equivariant average
\[
 J_\nu F(g)=
 \int \psi_\nu(n)\overline{\psi'_\nu(n')}\,
 \pi_\nu(w)F(w^{-1}n^{-1}gn'w')\pi'_\nu(w'^{-1})
 \,dw\,dn\,dw'\,dn'.
                                                               \tag{AG.3}
\]
The displayed integral uses the source's fixed Haar measures
and the required convergence hypotheses.  In its big-cell coordinates
\((z_1,\ldots,z_r,z)\), where \(r=r_S\), the energy is
\[
 \mathcal J=\kappa^{-(r+1)}
 \prod_\nu\int_{\mathbb C^{r+1}}
 \left[\prod_{a=1}^{r}|z_a|^{\,2(2(r-a)+k+1)}\right]
 (|z|^2-T_uT_v)_+^{\,k-1}
 \|J_\nu F_\nu(z_1,\ldots,z_r,z)\|_{\rm HS}^2\,dz .
                                                               \tag{AG.4}
\]
Here Lebesgue measure is used at each complex coordinate, and
\(\kappa>0\) is the common covolume of the scaled arithmetic
lattices in \(\mathbb C^f\); the source chooses the block-dependent
scalings so these covolumes agree.  Boundary values of the
positive part are immaterial to the integral.  The scalar pair identity rewrites
\(\gamma_G\operatorname{Re}\sum_{\rm scalar}p\) as
\((n_u/\dim\pi_u)\) times a positive quadratic form in tensor
products of these amplitudes.  Its coefficient kernel is bounded,
in the positive-semidefinite order, by an arithmetic lattice
kernel.  This is the precise input to the next argument.

The words indexing \(u\) are binary.  A differing letter at position
\(\ell\) gives, at every \(\nu\), the two compact parameters
\[
 v_{\nu\ell}-\tfrac12-\Delta,\qquad v_{\nu\ell}-\tfrac12 .
\]
Thus
\[
 I_\nu=[\,\tfrac52-v_{\nu\ell},\ \tfrac{\Delta}{2}-v_{\nu\ell}\,]
                                                               \tag{AG.5}
\]
lies strictly between their negatives and stays at least \(2\)
from both endpoints when \(\Delta\ge5\).  Only one differing
letter is needed; this step uses no lower bound for Hamming distance.
The large width of \(I_\nu\), together with the remaining
Harish--Chandra gaps, permits repeated spectral differentiation.
The normalization of a differential word \(w\) divides by
\(\prod_{\nu,Y\in w_\nu}|P_{\lambda_\nu}(Y)|\), with
\(P_\lambda(Y)=\prod_j(Y+\lambda_j)\).

Here is the local analytic input that still requires the
representation-theoretic proof.  At each level its Fourier transform
has zero mode equal to the preceding level's energy and satisfies
\[
 |\widehat{\mathcal N}_w(\xi)|
 \le\prod_\nu\min\{1,(x_{i\nu}/|\xi_\nu|)^K\}
             \sum_{w'}\theta_\xi(w')\mathcal E_{w'},\qquad
 \theta_\xi\ge0,\quad\sum_{w'}\theta_\xi(w')=1.                 \tag{AG.6}
\]
The extensions \(w'\) use at most \(K\) additional letters.
A factor with \(\xi_\nu=0\) is interpreted as \(1\).
There are corresponding off-diagonal estimates with
\(\xi_\nu=t_\nu^{-1}-(t'_\nu)^{-1}\), and estimates for blocks of
size \(s\ge2\).  The source takes \(K=5\).  These estimates,
positivity of the intermediate forms, and their common limiting
energy are substantive inputs, not definitions of a proof.

\paragraph{Why the error has a power \(f\) (AG6--AG7).}
Let \(E\) be the auxiliary CM field of degree \(2f\), with exactly
one embedding from each conjugate pair.  If \(\mathfrak a\ne0\)
is a fractional ideal and
\(\prod_\nu x_\nu\le \sqrt{N\mathfrak a}/\Theta^f\), the elementary
packing argument in the appendix proves
\[
 \sum_{0\ne z\in\mathfrak a}
       \prod_\nu\min\{1,(x_\nu/|\nu z|)^K\}
 \le \bigl(384\,\Theta^{-(K-1/2)}\bigr)^f
 \quad(K\ge3,\ \Theta\ge1).                                  \tag{AG.7}
\]
This is a new sufficient replacement for the weaker exponent
\(K-5/2\) used in the archive.  No improvement to the final RH
exponent is claimed.  The source defines \(\Theta\) from the
maximum of all relevant spectral-size products divided by square
roots of fractional-ideal norms, so the same \(\Theta\) controls
both Fourier and off-diagonal sums.

Poisson summation now separates the zero frequency from the
nonzero frequencies in (AG.6).  The former contributes exactly the
preceding energy; the latter contributes the small quantity in
(AG.7).  Splitting the final block of a positive Gram form and
using its triangle inequality gives the following concrete
recurrence for square roots of level energies:
\[
 x_i\le(1+\varepsilon)\sqrt{1+e_i}\,x_{i-1}
       +\sum_{s=2}^{i}a^{s-1}x_{i-s},\qquad
 a=(C/\Theta)^f,\quad 0\le\varepsilon\le a.                   \tag{AG.8}
\]
The \(e_i\ge0\) are the arithmetic kernel's density losses.
For the finitely many place patterns, the source closes the
finite supports of these probability distributions under the
bounded number of permitted extensions.  Taking a maximum over
that finite set bounds every average without a factor equal to
the number of words.  Finiteness follows because every nontrivial
extension strictly increases an integer total length, while the
length budget is bounded; the resulting tree has finite depth
and finite branching.

For completeness, suppose \(a\le1/2\) and \(na\le1\), where
\(n=r+1\).  Divide \(x_i\) by
\(\prod_{j\le i}\sqrt{1+e_j}\) and call the result \(y_i\).
Since the removed factors are at least one,
\[
 y_i\le(1+\varepsilon)y_{i-1}+
              \sum_{s=2}^{i}a^{s-1}y_{i-s}.
\]
If \(y_0=0\), induction gives \(y_i=0\) for all \(i\).
Otherwise induction gives \(y_i\le(1+3a)^iy_0\): the terms with \(s\ge2\),
after division by \((1+3a)^{i-1}y_0\), sum to at most
\(a/(1-a)\le2a\).  Consequently
\[
 x_n^2\le
 [1+(e^6-1)na]\prod_{j=1}^{n}(1+e_j)\,x_0^2.                \tag{AG.9}
\]
Indeed \((1+3a)^{2n}\le e^{6na}\le1+(e^6-1)na\).
This proves the leading constant \(1\) and the stated small
relative error.  Increasing the fixed constant \(C\) absorbs
the harmless absolute constants.  Taking finite sums first,
then the smooth-cutoff limit and increasing finite sets, is
justified by the stated positivity and absolute summability.

\paragraph{Returning to the period maps (AG8--AG10).}
The mass and formal-degree normalization, a separate arithmetic
input, bounds the base energy after multiplication by
\(n_uc_W/\dim\pi_u\) by \(e^{8q_S}X_{uv}\), where \(c_W>0\)
is the scalar in the arithmetic weight identity and
\(q_S=\sum_{\mathfrak p}(N\mathfrak p)^{-k}\).
The density factors and \(e^{8q_S}\) are absorbed into
\(1+e_{\rm loc}\).  Equations (AG.6)--(AG.9) therefore yield
\[
 S_{uv}\le
 [1+C(r+1)(C/\Theta)^f](1+e_{\rm loc})F X_{uv}
 \quad(u\ne v).                                             \tag{AG.10}
\]
The source's parameter comparison supplies
\((r+1)(C/\Theta)^f\le b^{-7}\), with \(b=e^L\).
It also gives, for \(\delta=b^{-7}\),
\[
 F\le1+\delta,\quad E\le\delta,\quad e_{\rm loc}\le\delta,
 \quad \|N_{uv}\|_{\rm HS}^2\le\delta^3X_{uv}.
\]
Use
\(\|N+G\|^2\le(1+\delta)\|G\|^2+
(1+\delta^{-1})\|N\|^2\) in (AG.1)--(AG.2).
The result is at most
\([1+(13+8C)\delta]X_{uv}\), hence
\[
 \|K_{uv}\|_{\rm HS}^2\le(1+b^{-6})\,n_un_v/D_*
\]
once \(b\ge13+8C\).
On the diagonal the scalar sum needs only a coarse
\(e^{C_E m^5fL}n_u^2/D_*\) bound.  The same decomposition gives
\[
 \|K_{uu}-Z_{uu}-I\|_{\rm HS}^2\le\Gamma n_u^2/D_*,
 \qquad
 \|Z_{uu}\|_{\rm HS}^2\le n_u e^{-kf\log k}.
\]
Since \(n_u\Gamma^2/D_*\le C_0\), it follows that
\[
 \|K_{uu}-I\|_{\rm HS}^2\le
 \bigl(2C_0/\Gamma+2e^{-kf\log k}\bigr)n_u .                 \tag{AG.11}
\]
A separate automorphic dimension estimate supplies
\[
 D_{\rm all}\le(1+\eta_1)D_*,
 \qquad \eta_1=e^{-c_7Rf\log R},\quad c_7>0,
\]
where \(R\) is the source's total rank, not \(r=r_S\).
This counting estimate is a retained input; it does not follow
from the period-kernel calculation.
These inequalities give the Gram hypotheses used below.
The retained obligations are now more specific:
the geometric expansion and trace identity, the local spectral
identities underlying (AG.6), the arithmetic kernel and mass
normalization, and the loop, nongeneric defect-one, and
zero-determinant diagonal estimates.
